\section{Complex powers of a cotangent resolvent}
\label{power:sec}

The rational resolvent theorem in the incoming harmonic-parity report
leaves open its extension to complex powers (Question R4).  The
complex-power sine transform in the higher-reflection report treats a
different, reflection-symmetric weight.  We now give a one-sided
complex-power transform at every nonreal resolvent parameter, an ordinary
Euler integral involving a single classical polylogarithm, and a completed
germ at every integral power.  The latter includes a removable but
nonzero endpoint correction when the integral power is negative.
Classical Hurwitz Fourier analysis and beta integrals supply the analytic
tools \cite{DLMF,EspinosaMoll}; the formulas and their
normalizations below are proved directly.

\subsection{A branch fixed by the interval}

For $z\in\mathbb C\setminus\mathbb R$, put
\begin{equation}\label{power:coordinates}
 \epsilon=\operatorname{sgn}\Im z,\quad
 a=z+\epsilon\pi i,\quad q=\frac{z-\epsilon\pi i}{a},\quad
 \sigma=-2\epsilon\pi i,\quad C=-\frac1a=\frac{1-q}{\sigma}.
\end{equation}
Thus $|q|<1$.  Throughout,
\[
 \Log\sigma=\log(2\pi)-\epsilon\pi i/2,
 \qquad \Log C=\Log(1-q)-\Log\sigma.
\]
Both $\Log C$ and $\Log(1-q)$ in these equations are principal
logarithms.  Write
\begin{equation}\label{power:R}
 K(x)=\pi\cot(\pi x),\quad R_z(x)=\frac1{K(x)-z},\quad
 w=e^{2\epsilon\pi ix}.
\end{equation}
Direct algebra gives
\begin{equation}\label{power:factor}
 R_z(x)=C\frac{1-w}{1-qw}.
\end{equation}
The continuous logarithm of $R_z(x)$ on $0<x<1$ is fixed by
$\Log R_z(x)-\log x\to0$ at zero.  Its imaginary part lies between
$0$ and $\epsilon\pi$, and it approaches $\epsilon\pi$ at one.
In particular,
\begin{equation}\label{power:branches}
 R_z(x)^\lambda=\exp(\lambda\Log R_z(x)),\qquad
 R_z(1-y)^\lambda=e^{\epsilon\pi i\lambda}y^\lambda
                       V_z(-y)^\lambda,
 \quad V_z(x)=\frac{R_z(x)}x.
\end{equation}
Here $V_z$ is analytic near zero and $V_z(0)=1$.  Powers of $V_z$
use the analytic logarithm vanishing at zero.  These prescriptions also
specify all logarithmic derivatives in $\lambda$.

\subsection{A beta mixture of ordinary polylogarithms}

Define coefficients and their Dirichlet transform by
\begin{align}
 \left(\frac{1-w}{1-qw}\right)^\lambda
   &=\sum_{n\ge0}d_n(\lambda,q)w^n,\quad d_0=1,
       \label{power:dn}\\
 d_n(\lambda,q)&=\sum_{j=0}^n
       \frac{(-\lambda)_j(\lambda)_{n-j}}{j!(n-j)!}q^{n-j},
       \label{power:dnfinite}\\
 \mathcal L_\lambda(\nu,q)&=\sum_{n\ge1}
                               \frac{d_n(\lambda,q)}{n^\nu}.
       \label{power:L}
\end{align}
The last series is initially normally convergent for
$\Re(\nu+\lambda)>0$, on compact parameter sets away from the boundary
of $|q|<1$.  At integral $\lambda$ it can have a larger convergence
domain; our continuation will not rely on such exceptional improvements.
The notation in \eqref{power:L} is a transform of the explicit
coefficients \eqref{power:dnfinite}, not a claim that an unevaluated
integral has been reduced to familiar constants.

\begin{theorem}[The Euler--polylogarithm formula]\label{power:eulerthm}
For $-1<\Re\lambda<1$, $|q|<1$, and initially $\Re\nu>1$,
\begin{equation}\label{power:euler}
 \boxed{\displaystyle
 \mathcal L_\lambda(\nu,q)=
 -\frac{\sin\pi\lambda}{\pi}(1-q)
 \int_0^1 y^{-\lambda}(1-y)^\lambda
       \frac{\Li_\nu(q+(1-q)y)}{q+(1-q)y}\,\dd y.}
\end{equation}
The quotient $\Li_\nu(t)/t$ has value $1$ at $t=0$.  Thus the
formula remains valid when the segment from $q$ to $1$ passes through
zero.  Within the same $\lambda$ strip the formula also holds on
$\Re(\nu+\lambda)>0$ by analytic continuation of the convergent
integrals.  All fixed parameter derivatives commute with these integrals
on compact subsets of their convergence domains.
\end{theorem}

\begin{proof}
We first prove the coefficient identity
\begin{equation}\label{power:beta-coeff}
 d_n(\lambda,q)=-\frac{\sin\pi\lambda}{\pi}(1-q)
 \int_0^1y^{-\lambda}(1-y)^\lambda
                 (q+(1-q)y)^{n-1}\,\dd y,\qquad n\ge1.
\end{equation}
For $0<q<1$, deform the Cauchy coefficient contour for
$((1-w)/(1-qw))^\lambda$ onto $1<w<1/q$.  The difference of its
upper and lower boundary values is
$-2i\sin(\pi\lambda)((w-1)/(1-qw))^\lambda$.
The endpoint integrals exist precisely in the stated strip.  The
integral at infinity vanishes for $n\ge1$.  Substituting $t=1/w$
and then $t=q+(1-q)y$ proves \eqref{power:beta-coeff}.
Alternatively, its right side can be expanded as a finite beta sum
and checked against \eqref{power:dnfinite}.  Both sides are
polynomials in $q$, so the identity holds for all complex $q$.

For $\Re\nu>1$ we may sum \eqref{power:beta-coeff} after multiplying
by $n^{-\nu}$.  The segment lies in the closed unit disk and only
its last endpoint can meet the unit circle.  The absolutely convergent
polylogarithm series, together with the integrable beta weight, provides
domination.  This yields \eqref{power:euler}.  Near $y=1$ the possible
singular part of $\Li_\nu(q+(1-q)y)$ has order
$(1-y)^{\nu-1}$, with the usual logarithmic limits at integral
orders \cite[Eq.~25.12.12]{DLMF}.  The conditions
$\Re\lambda>-1$ and $\Re(\nu+\lambda)>0$ therefore give the
asserted extension and the same domination after any fixed number of
derivatives.  Near $y=0$, $\Re\lambda<1$ suffices, including $q=0$.
\end{proof}

There is a useful independent Mellin representation:
\begin{equation}\label{power:mellinL}
 \mathcal L_\lambda(\nu,q)=\frac1{\Gamma(\nu)}
 \int_0^\infty t^{\nu-1}
 \left\{\left(\frac{1-e^{-t}}{1-qe^{-t}}\right)^\lambda-1\right\}
 \,\dd t,
\end{equation}
initially for $\Re\nu>\max(0,-\Re\lambda)$.
At zero, the expression in braces is $t^\lambda$ times an analytic
function minus $1$; at infinity it decreases exponentially.
Finite Taylor subtraction in \eqref{power:mellinL} proves joint
meromorphic continuation in $(\nu,\lambda)$.  Its only possible
moving polar hyperplanes are
$\nu+\lambda=-j$ for $j\in\mathbb Z_{\ge0}$; some disappear
through zeros of $1/\Gamma(\nu)$ or of their coefficients.
This supplies a continuation independently of the Euler integral's
restricted $\lambda$ strip.

\subsection{The Hurwitz transform and an ordinary log-Gamma integral}

\begin{theorem}[Complex-power Hurwitz transform]\label{power:hurwitzthm}
For $p\in\mathbb Z_{\ge0}$, $\Re\lambda>-1$, and
$\Re(u+\lambda)>p$, away from the Hurwitz pole when $p=0$,
\begin{equation}\label{power:hurwitz}
 \boxed{\displaystyle
 \int_0^1\partial_x^p\zeta(1-u,x)R_z(x)^\lambda\,\dd x
 =C^\lambda\Gamma(u)\sigma^{p-u}
                    \mathcal L_\lambda(u-p,q).}
\end{equation}
The identity then continues meromorphically in its parameters.  Also,
for every $\Re\lambda>-1$,
\begin{equation}\label{power:mean}
 \int_0^1R_z(x)^\lambda\,\dd x=C^\lambda.
\end{equation}
\end{theorem}

\begin{proof}
For $\Re u>p+1$, the Hurwitz Fourier series and its first $p$
argument derivatives are absolutely convergent.  The nonzero Fourier
pairings are
\[
 \int_0^1\partial_x^p\zeta(1-u,x)e^{2\epsilon\pi inx}\,\dd x
     =\Gamma(u)\sigma^{p-u}n^{p-u},\qquad n\ge1,
\]
and the zero mode vanishes.  Multiply by the Abel-regularized expansion
\eqref{power:dn} and then let the Abel radius tend to one.
For $\Re\lambda>-1$ its boundary values are uniformly integrable:
near either endpoint they are bounded by a constant times
$x^{\min(\Re\lambda,0)}$ or $(1-x)^{\min(\Re\lambda,0)}$.
This proves the formula in an initial open region.  The Dirichlet
series is normally convergent under $\Re(u+\lambda)>p$.
The original integral has the same sufficient local conditions, since
its only nonsmooth Hurwitz term at zero is
$\prod_{j=1}^p(u-j)x^{u-p-1}$.  Holomorphic continuation within
that domain proves \eqref{power:hurwitz}.  The constant Fourier
coefficient of the Abel kernel is $C^\lambda$; the same integrable
limit proves \eqref{power:mean}.
\end{proof}

\begin{corollary}[An ordinary Gamma moment]\label{power:gamma}
For $\Re\lambda>-1$,
\begin{align}\label{power:gammaformula}
 \int_0^1\log\Gamma(x)R_z(x)^\lambda\,\dd x
 =C^\lambda\left\{\frac{\log(2\pi)}2
 +\frac{(\gamma+\Log\sigma)\mathcal L_\lambda(1,q)
       -\partial_\nu\mathcal L_\lambda(1,q)}{\sigma}\right\}.
\end{align}
\end{corollary}
\begin{proof}
Differentiate \eqref{power:hurwitz} with $p=0$ at $u=1$ and use
$\zeta'(0,x)=\log\Gamma(x)-\tfrac12\log(2\pi)$ and
$\Gamma'(1)=-\gamma$ \cite{DLMF}.
The logarithmic endpoint singularity is integrable under the stated
condition, so this is an ordinary integral and an ordinary derivative.
\end{proof}

Thus, in the strip of Theorem~\ref{power:eulerthm}, the Gamma moment
is an explicit combination of two integrals of $\Li_1$ and its first
order derivative against an elementary beta weight.  At $q=0$, its
first transform is $\mathcal L_\lambda(1,0)=-\psi(1+\lambda)-\gamma$.
Section~\ref{moment:sec} gives the stronger finite Gamma generating
formula for all polynomial moments.

\subsection{Completion at every integral power}

Let
\begin{equation}\label{power:stieltjes}
 g(u,x)=\zeta(1-u,x)+\frac1u
     =\sum_{m\ge0}\gamma_m(x)\frac{u^m}{m!},\qquad
 A_p(u)=\prod_{j=1}^p(u-j),\quad A_0=1.
\end{equation}
Define finite local coefficients
\begin{align}
 v_j(\lambda;z)&=[x^j]V_z(x)^\lambda,\label{power:vj}\\
 h_p(u,x)&=\partial_x^p\left\{\zeta(1-u,1+x)+\frac1u\right\},
        \label{power:hp}\\
 b_j(u,\lambda;z)&=[x^j]h_p(u,x)V_z(x)^\lambda.
        \label{power:bj}
\end{align}
The function $h_p$ is holomorphic near $(u,x)=(0,0)$.  Its coefficients
are explicit: the coefficient of $x^j$ is
\begin{equation}\label{power:hcoeff}
 \frac{(-1)^{p+j}(1-u)_{p+j}}{j!}\zeta(1-u+p+j)
 \quad(p+j>0),
\end{equation}
and is $\zeta(1-u)+1/u$ when $p=j=0$.
The first coefficients of $V_z$ are
\[
 V_z(x)=1+zx+(z^2+\pi^2/3)x^2
                   +(z^3+2\pi^2z/3)x^3+O(x^4).
\]
All fixed Taylor jets of $v_j$ and $b_j$ required for coefficient
extraction are therefore finite polynomials in $z$, the expansion
parameters, zeta values or their derivatives, and Stieltjes
coefficients.  No new global integral enters them.

For an integer $L$, set $\lambda=L+v$ and
\begin{align}
 \mathcal T_p(u,\lambda;z)
   &=C^\lambda\Gamma(u)\sigma^{p-u}
                    \mathcal L_\lambda(u-p,q)
         +\mathbf1_{p=0}\frac{C^\lambda}{u},\label{power:T}\\
 \mathcal H_{p,L}(u,v;z)
   &=\mathcal T_p(u,L+v;z)
      -\mathbf1_{L\le p}
        \frac{A_p(u)v_{p-L}(L+v;z)}{u+v}\notag\\
   &\quad-\mathbf1_{L\le-1}
       \frac{\{1+(-1)^{-L-1}e^{\epsilon\pi i(L+v)}\}
                    b_{-L-1}(u,L+v;z)}{v}.
       \label{power:H}
\end{align}
A term with a false indicator is omitted completely.  In the last
line the quotient is read by its removable analytic extension.

\begin{theorem}[Joint Stieltjes and power completion]\label{power:completion}
For each $p\ge0$, $L\in\mathbb Z$, and $z\notin\mathbb R$,
$\mathcal H_{p,L}$ is holomorphic near $(u,v)=(0,0)$.
For all integers $m,k\ge0$,
\begin{equation}\label{power:allmoments}
 \boxed{\displaystyle
 \FP\int_0^1\gamma_m^{(p)}(x)R_z(x)^L
                    (\Log R_z(x))^k\,\dd x
       =m!k![u^mv^k]\mathcal H_{p,L}(u,v;z).}
\end{equation}
The finite part is taken in the fixed coordinates $x$ and $1-x$.
All coefficient extractions and all $z$ derivatives are locally
holomorphic on each open half-plane.
\end{theorem}

\begin{proof}
The exact local splitting is
\[
 \partial_x^p g(u,x)=A_p(u)x^{u-p-1}+h_p(u,x).
\]
At zero, multiplication by $R_z(x)^\lambda=x^\lambda V_z(x)^\lambda$
produces the two families
\[
 A_p(u)v_j(\lambda;z)x^{u+\lambda-p-1+j},\qquad
 b_j(u,\lambda;z)x^{\lambda+j}.
\]
After setting $\lambda=L+v$, their integration denominators pass
through zero only for $j=p-L\ge0$ and $j=-L-1\ge0$, respectively.
They give the second line of \eqref{power:H} and the first endpoint's
contribution to its third line.

At the other endpoint put $y=1-x$.  The smooth Hurwitz factor is
exactly $h_p(u,-y)$ and \eqref{power:branches} applies.  Its
coefficient of $y^j$ is
$e^{\epsilon\pi i\lambda}(-1)^jb_j(u,\lambda;z)$.
This supplies precisely the remaining term in the braces in
\eqref{power:H}.  Their sum vanishes on $v=0$, since
$(-1)^{-L-1}e^{\epsilon\pi iL}=-1$.  Nevertheless its quotient by
$v$ must be subtracted with the complete numerator.

To make the argument analytic, choose a small endpoint interval and
subtract finitely many of the displayed monomials, enough to leave
a uniformly integrable remainder throughout a small parameter
polydisk.  The middle interval is holomorphic.  Nonresonant local
denominators are bounded away from zero.  For a resonant monomial
with exponent $\alpha-1$, the difference between its meromorphic
integral and the generating function of its coordinate finite parts
is its full amplitude divided by $\alpha$, because
\[
 \frac{\delta^\alpha}{\alpha}
 =\frac1\alpha+\frac{\delta^\alpha-1}{\alpha}.
\]
The second summand is entire in $\alpha$ and its derivatives give
the constants of the corresponding logarithmic integrals.  This
also shows that the arbitrary splitting point cancels out.  Finally,
the global Hurwitz pole contributes $-C^\lambda/u$ only for $p=0$,
and \eqref{power:T} has already removed it by \eqref{power:mean}.
The remaining germ is holomorphic and its coefficients are exactly
\eqref{power:allmoments}.  The same normally convergent construction
is uniform on compact subsets of either $z$ half-plane.
\end{proof}

Away from integral powers there is no resonant exponent at $u=0$.
Consequently, for $\lambda\notin\mathbb Z$, the same local subtraction
proof gives the simpler identity
\begin{equation}\label{power:noninteger}
 \FP\int_0^1\gamma_m^{(p)}(x)R_z(x)^\lambda
             (\Log R_z(x))^k\,\dd x
 =m![u^m]\partial_\lambda^k\mathcal T_p(u,\lambda;z).
\end{equation}
Here the finite part means the constant term in the generalized
power--logarithm endpoint expansion.  For nonintegral $\lambda$,
none of its subtracted endpoint powers has exponent zero after
integration, so differentiation and this finite part commute locally
in $\lambda$.  At integral powers the full completion
\eqref{power:H}, rather than a limit of \eqref{power:noninteger},
specifies the coordinate finite part.

\paragraph{The correction at negative powers is not zero.}
Although the last quotient in \eqref{power:H} is removable, its
value at $v=0$ is
\begin{equation}\label{power:smoothcontact}
 -\epsilon\pi i\,b_{-L-1}(u,L;z),\qquad L\le-1.
\end{equation}
Removing only pole residues would miss it.  It records the different
endpoint phases of $R_z^\lambda$ before $\lambda$ is specialized to
an integer.  At a fixed negative integer the integrand itself is a
polynomial in $K(x)-z$, so the resulting finite parts must have the
corresponding polynomial dependence on $z$.  Formula
\eqref{power:smoothcontact} is necessary for that compatibility.

The same issue already appears without a Hurwitz factor.  For $M\ge1$,
\begin{equation}\label{power:negative-mean}
 \FP\int_0^1(K(x)-z)^M\,\dd x
 =\frac{(-z+i\pi)^M+(-z-i\pi)^M}{2},
\end{equation}
whereas the continued mean in \eqref{power:mean} is
$C^{-M}=(-z-\epsilon i\pi)^M$.  Reflection gives zero for odd
cotangent powers, and integration of their derivative recurrence gives
$\FP\int_0^1K^{2j}\,\dd x=(-\pi^2)^j$, proving
\eqref{power:negative-mean}.  The difference is exactly
$\epsilon\pi i\,[x^{M-1}]V_z(x)^{-M}$, as the two-endpoint
completion predicts.  This is an elementary check at every negative
integral power.
An independent all-$M$ residue proof is given in
Appendix~\ref{contact:verification}.

\begin{example}[Two normalization checks]\label{power:canaries}
For $L=0$ and $k=0$, the theorem gives
$\FP\int_0^1\gamma_m^{(p)}(x)\,\dd x=0$.
For $L=-1$, $p=m=k=0$, it gives
\begin{equation}\label{power:cotcheck}
 \FP\int_0^1\gamma_0(x)(\pi\cot\pi x-z)\,\dd x
        =-\frac{\pi^2}{2}.
\end{equation}
Indeed $\mathcal L_{-1}(u,q)=(1-q)\zeta(u)$,
$v_1(-1;z)=-z$, and $b_0=\zeta(1-u)+1/u$.
At $v=0$, \eqref{power:H} reduces, by the zeta functional equation, to
\[
 2\pi\Gamma(u)(2\pi)^{-u}\sin(\pi u/2)\zeta(u),
\]
whose value at zero is $-\pi^2/2$.  The same value follows directly
from digamma reflection and $\FP\int_0^1K(x)^2\,\dd x=-\pi^2$.
For $L=1$ and $p=1$, the moving term is $(u-1)/(u+v)$; at
$z=i\pi$ it recovers
\[
 \FP\int_0^1\frac{\psi'(x)}{\pi\cot\pi x-i\pi}\,\dd x
       =1-\gamma-\log(2\pi)+i\pi/2,
\]
including the inherited resolvent theorem's harmonic correction.
\end{example}

\subsection{An explicit logarithmic resolvent identity}

At first order in $\lambda$ there is no need to retain the transform
$\mathcal L$: differentiating \eqref{power:dn} at zero gives
\begin{equation}\label{power:firstlambda}
 \left.\partial_\lambda\mathcal L_\lambda(u,q)\right|_{\lambda=0}
       =\Li_{1+u}(q)-\zeta(1+u).
\end{equation}

\begin{corollary}[All Stieltjes indices with one logarithm]
\label{power:logcor}
The expression in braces below is holomorphic at $u=0$, and
\begin{equation}\label{power:alllog}
 \boxed{\displaystyle
 \FP\int_0^1\gamma_m(x)\Log R_z(x)\,\dd x
 =m![u^m]\left\{
 \Gamma(u)\sigma^{-u}\bigl(\Li_{1+u}(q)-\zeta(1+u)\bigr)
       +\frac{\Log C}{u}+\frac1{u^2}\right\}.}
\end{equation}
In particular, if a superscript $[1]$ denotes differentiation of the
polylogarithm with respect to its order, then
\begin{align}\label{power:logexplicit}
 \FP\int_0^1\gamma_0(x)\Log R_z(x)\,\dd x
  ={}&\gamma_1+\frac{\gamma^2}{2}-\frac{\zeta(2)}2
              -\frac{(\Log\sigma)^2}{2}\notag\\
    &+\Li_1^{[1]}(q)
               +(\gamma+\Log\sigma)\Log(1-q).
\end{align}
\end{corollary}
\begin{proof}
Differentiate \eqref{power:H} at $v=0$ with $p=L=0$, using
\eqref{power:firstlambda}.  The derivative of $-1/(u+v)$ is
$1/u^2$.  Expansion of
\[
 \Gamma(u)\sigma^{-u}
 =\frac1u\exp\left(-Au+\sum_{j\ge2}\frac{(-1)^j\zeta(j)}j u^j\right),
 \qquad A=\gamma+\Log\sigma,
\]
and $\zeta(1+u)=u^{-1}+\gamma-\gamma_1u+O(u^2)$ proves
\eqref{power:logexplicit}.  The coefficient of $u^{-1}$ cancels
because $\Li_1(q)=-\Log(1-q)$ and
$\Log C=\Log(1-q)-\Log\sigma$.
\end{proof}

At $z=\epsilon i\pi$, the polylogarithm terms disappear and the
identity becomes the compact evaluation
\begin{equation}\label{power:logspecial}
 \FP\int_0^1\gamma_0(x)
       \left\{\log\frac{\sin\pi x}{\pi}+\epsilon\pi ix\right\}\,\dd x
 =\gamma_1+\frac{\gamma^2}{2}+\frac{\pi^2}{24}
       -\frac{\log^2(2\pi)}2+\frac{\epsilon\pi i\log(2\pi)}2.
\end{equation}
The imaginary part is also an ordinary consequence of Raabe's
log-Gamma integral.  The constant real part supplies a useful
independent check of the branch and finite-part conventions.
The $q=0$ specialization is compatible with the prior sine-power
completion; the arbitrary nonreal $z$ identity is the new extension.

\subsection{Exact derivatives and normalized primitives in the resolvent parameter}

Denote the left side of \eqref{power:allmoments} by
$M_{m,p;L,k}(z)$.  Coefficientwise differentiation in $z$ is legitimate
by Theorem~\ref{power:completion}.  Since
$\partial_z\Log R_z=R_z$, we obtain the exact ladder
\begin{equation}\label{power:zladder}
 \frac{\dd}{\dd z}M_{m,p;L,k}(z)
    =L\,M_{m,p;L+1,k}(z)+k\,M_{m,p;L+1,k-1}(z),
\end{equation}
with the last term omitted for $k=0$.  This differentiates the actual
fixed-coordinate moments, rather than an uncompleted spectral limit.

For $L\ne1$ define
\begin{equation}\label{power:primitive}
 P_{m,p;L,k}(z)=
 \sum_{j=0}^k\frac{(-1)^jk!}{(k-j)!(L-1)^{j+1}}
                         M_{m,p;L-1,k-j}(z).
\end{equation}
For $L=1$ put $P_{m,p;1,k}=M_{m,p;0,k+1}/(k+1)$.
Equation \eqref{power:zladder} telescopes to $P'=M$.
Consequently
\begin{equation}\label{power:normalizedprimitive}
 \int_{z_0}^{z}M_{m,p;L,k}(\xi)\,\dd\xi
       =P_{m,p;L,k}(z)-P_{m,p;L,k}(z_0)
\end{equation}
for any path within the same open half-plane.  This fixes every additive
constant explicitly.  In particular \eqref{power:logexplicit} is an
explicit primitive of the ordinary $\gamma_0$ resolvent integral.

\paragraph{An exact boundary counterexample to naive substitution.}
The meromorphic transform is not, at every exceptional point, the
pointwise integral of the specialized integrand.  At $p=2$, $u=1$,
$\lambda=1$, the latter is zero because
$\partial_x^2\zeta(0,x)=0$.  Continuing \eqref{power:hurwitz} along
$\lambda=1$ first gives $-1$.  Indeed
$\mathcal L_1(u-2,q)=(q-1)\Li_{u-2}(q)/q$, with its removable
$q=0$ interpretation, gives $-1$ at $u=1$ by elementary algebra.
The sufficient convergence domain has $\Re(u+\lambda)>p$ and
excludes this boundary.  Locally the discrepancy is the limit of
$(u-1)(u-2)/(u-1)=u-2$.  This example identifies a concrete
normalization error in an uncompleted extension; it is not an allegation
that the inspected incoming resolvent theorem makes that error.
